\section{A candidate proof of uniform physical-word pressure}
\label{sec:wp-pressure}
\textbf{Author working draft for independent reconstruction.}
This file is not a published assertion that the pressure input is
discharged.  The complete-barrier and small-cell arguments below are the
new geometric steps to be checked before integration.  No existing
proof body or historical claim is replaced.

The system, unit-speed clock, Euclidean triangular lattice and radius
interval $R\in[9/20,47/100]$ are unchanged.  The pressure sum is precisely
\[
 P_m(R)=\sum_{w\in\mathcal W_m(R)}
       \sup_{x\in[w]_R\cap M_R'}J^sT_R^m(x),\qquad P_0=1,
\]
where a word records every actual lifted collision.  The supremum is
over the set regular in both infinite directions, and stable arclength
means Euclidean arclength in $(s,\varphi)$.
The proof uses no conditioning or exact-label lower bound.

\subsection{Constants and the actual backward partition}

Write
\[
 a=100/47,\quad K=170/9,\quad H=19,\quad
 v=95/9,\quad \Gamma=\sqrt{1+(20/9)^2},\quad
 \Lambda=59/47,\quad e_0=2a/H^2.
\]
The fixed line of slope $-v$ belongs to every physical stable cone.
Let $\mathcal H$ be the original $q=2$ homogeneity partition, with
$H_k^\pm$ having angular width at most $3k^{-3}$.
Increase the fixed threshold $k_0$, if necessary, so that the actual
one-flight count $J_f$ satisfies
\[
 \theta:=\Lambda^{-1}
       +\frac{J_f}{a\tau_-(k_0-1)}<1,\qquad \tau_-=3/50.
\]
The existing curvature and curve-growth theorems remain valid after
this fixed refinement and fixed shortening of the admissible length.
No parameter in this construction depends on the collision count.

Let $\mathcal B_n(R)$, for $n\ge1$, be the connected homogeneous
backward cells obtained by retaining all genuine past singularities
and the homogeneity cuts from times $-n,\ldots,0$.
In the notation of Baladi--Demers (2.4), this is
$\mathcal M_{-n}^{0,H}\vee\mathcal H$.
It determines the $n$ actual past physical symbols.
Set $\mathcal B_0=\mathcal H$.
For $y\in M_R'$ put
\[
 k_n(y)=J^uT_R^{-n}(y),\qquad
 r_n(y)=c(T_R^{-n}y)k_n(y),\qquad
 A_n(B)=\sup_{B\cap M_R'}r_n,\qquad
 S_n=\sum_{B\in\mathcal B_n}A_n(B).
\]
Thus $k_0=1$ in this display only; the homogeneity threshold is always
written as a subscript on $H_k$.  To avoid this notational collision
in the final manuscript, rename the empty cocycle $\mathfrak k_0$.

The exact area identity in module 313 gives, with
$x=T_R^{-n}y$,
\begin{equation}\label{eq:wp-area}
 r_n(y)=c(y)\,j_n(x)\frac{e(y)}{e(x)},\qquad
       e_0\le e\le1.
\end{equation}
This last incidence is essential: it will pay the countably many
homogeneity strips in the transverse tests.
Also
\[
 r_{l+j}(y)=
 r_l(T_R^{-j}y)\,k_j(y).
\]
This is an exact physical cocycle identity.

For each nonzero possible previous-center displacement $z$, use the
collision chart
\[
 \alpha-\arg z\in(-\pi/2,\pi/2),\qquad
 s_z=R(\alpha-\arg z),\qquad
 h_z=\varphi+v s_z.
\]
Every regular state whose previous disk has center $z$ is in this chart.
Indeed, if its previous contact is $q'=z+R n'$ and its current contact
is $q=Rn$, then
\[
 z\cdot n=R(1-n'\cdot n)+(q'-q)\cdot n
         =R(1-n'\cdot n)+\tau c>0.
\]
There are at most $N_f$ choices of $z$ and the chart width is $\pi R$,
so no artificial seam is inserted at each iterate.
Each $B\in\mathcal B_n$, $n\ge1$, uses the chart specified by its
last past symbol.

\subsection{The barrier statement requiring geometric reconstruction}

\begin{lemma}[Complete barriers and transverse fibres]
\label{lem:wp-barriers}
In the above chart a backward cell is a band in the coordinates
$(h_z,s_z)$:
\[
 B=\{(h,s):h\in I_B,\ g_B^-(h)<s<g_B^+(h)\},
\]
where $I_B$ is an interval and both boundary functions are
nondecreasing and $1/v$-Lipschitz.
The statement includes limits of countably many homogeneous barriers.
A stable line $h=\mathrm{constant}$ therefore meets $B$ in an interval.
Moreover
\begin{equation}\label{eq:wp-fibre}
 \sup_h(g_B^+(h)-g_B^-(h))
 \le C_s\Lambda^{-n},\qquad
 C_s=\Gamma\,\frac{\pi H}{a}.
\end{equation}
Consequently, if $|I_B|<\delta$, then
\begin{equation}\label{eq:wp-diameter}
 \operatorname{diam}(B)
 \le(2+v^{-1})\delta+(1+v)C_s\Lambda^{-n}.
\end{equation}
\end{lemma}

\begin{proof}[Candidate proof; endpoints and complete-barrier extension to be checked]
Each current homogeneity boundary is horizontal.  Every regular
positive iterate of a past true singularity or past homogeneity
boundary is a graph of slope in $[a,K]$.  This follows from the
physical differential: a horizontal starting tangent has image
slope $\varkappa'+c'\varkappa/(\tau\varkappa+c)$, while an already
positive cone tangent remains in that cone.
The endpoints of an image arc lie on an older image singularity or
on the boundary of the collision cylinder.  An arc cannot terminate
freely in a regular interior.

Order the arcs by the first time at which they are introduced.
At either endpoint continue an arc along the older barrier it meets,
in the direction of increasing, respectively decreasing, $h$.
Each switch reduces the introduction time, so this extension ends
on the boundary and produces a complete graph over $h$ without adding
an interior cut.  Boundary segments of the fixed physical chart may
be included in the extension.
In $(h,s)$ an arc of slope $V\ge0$ has derivative
$ds/dh=(v+V)^{-1}\in[0,1/v]$; a vertical chart edge has derivative zero.
Thus every extended barrier is a nondecreasing $1/v$-Lipschitz graph.

A complementary cell has a fixed side relative to every complete
barrier.  At each $h$ its fibre is the interval between the supremum
of its lower barriers and the infimum of its upper barriers.
Their extrema are again nondecreasing $1/v$-Lipschitz functions.
The set of $h$ where the interval is nonempty can have several
components when barriers cross; taking a connected cell selects
one such interval.  This argument does not assert that two
increasing graphs cross only once.
For countably many strip barriers, exhaust their indices.  They
accumulate only on the genuine image singularities already present.
Equivalently, take the same supremum and infimum over all barriers.
No new regular open fibre is discarded.

A stable fibre in $B$ has one inverse branch through time $-n$.
Its inverse image is a cone-stable graph; its length is at most
$\pi H/a$, because its angle variation is at most $\pi$ and
$|d\varphi/ds|\ge a$.
The adapted expansion, with one endpoint metric comparison, is
at least $\Gamma^{-1}\Lambda^n$.
This proves \eqref{eq:wp-fibre}, with the harmless larger constant
displayed there.  If $h,h'\in I_B$, the Lipschitz boundary bounds give
$|s-s'|\le v^{-1}|h-h'|+C_s\Lambda^{-n}$.
Since $\varphi=h-vs$, the sum of the two coordinate differences gives
\eqref{eq:wp-diameter}.
\end{proof}

\begin{lemma}[Uniform small-cell one-step contraction]
\label{lem:wp-small}
There exists $\eta>0$, independent of $R$, such that every homogeneous
backward cell $B$ of Euclidean diameter at most $\eta$ satisfies
\begin{equation}\label{eq:wp-cell-step}
 \sum_{B'\subset T_RB}
             \sup_{B'\cap M_R'}J^{u,*}T_R^{-1}\le\theta<1.
\end{equation}
The sum is over the complete genuine one-flight/homogeneous
decomposition.  The same conclusion holds after intersecting $B$
with a sufficiently small coordinate rectangle; its vertical and
horizontal sides become positive cone barriers after one flight.
The star denotes the adapted cone norm of module 282.
\end{lemma}

\begin{proof}[Candidate two-dimensional first-hit proof]
Fix a limiting radius and outgoing state, allowing initial grazing.
Follow its actual oriented ray to its first strict disk intersection.
The strong finite-horizon argument supplies this disk; there are
only finitely many earlier candidate contacts, all tangencies.
The strict entry persists in a neighborhood.
Any branch hitting an earlier tangent disk has terminal incidence
tending to zero and therefore cannot have its image in $H_0$ after
the neighborhood is reduced.
A disk later than the persistent strict disk cannot be the first hit.
The positive feet of the finite candidate contacts are distinct,
and their first-entry order persists after this reduction.

Connectivity must be proved in two dimensions.
Use the oriented-line coordinates
\[
 \vartheta=s/R+\varphi,\qquad b=-R\sin\varphi.
\]
This is a homeomorphism on the outgoing collision cylinder in a
local lift, including its grazing boundary; it is a smooth
coordinate change in the regular interior.
For a candidate center $z$, its two tangencies are exactly
\[
 b=z\cdot v^\perp\pm R.
\]
A fixed terminal homogeneity strip replaces the two constants by
the corresponding two numbers in $[-R,R]$.
Every earlier near tangency imposes one fixed side of one such
graph.  The positive-foot condition already persists.
For distinct centers, the difference of any two such graph
functions is $d\cdot v^\perp-a'$ with
$|d|\ge1$ and $|a'|\le2R_+<1$.
Its consecutive zeros are separated by at least
$2\arccos(2R_+)>1/10$.
Restrict the angular neighborhood to width less than $1/10$.
Each comparison of an upper and lower candidate graph therefore
imposes an interval or half-interval in $\vartheta$, never two
disconnected allowed intervals.

A backward barrier, written as $b=\beta(\vartheta)$, has
\[
 \frac{d\beta}{d\vartheta}
       =-\frac{RcV}{\varkappa+V}\in[-Rc,0]
\]
for $V\ge0$, with the same bounds for horizontal and vertical sides.
At an intersection with a candidate graph its derivative difference is
\[
 -z\cdot v-\beta'(\vartheta)
       \le-z\cdot v+Rc=-t_z<0.
\]
The positive foot is bounded away from zero at these candidate
contacts; for a disk-intersection boundary one may use
$t_z\ge\sqrt{1-2R_+}>0$.
Thus a candidate graph crosses a backward boundary only in one
direction.  Each resulting inequality again imposes an interval
in $\vartheta$.
The barrier-band description and these inequalities show that
each selected first-hit center and target strip gives at most one
connected piece in the neighborhood.
In particular there is at most one $H_0$ piece, and at most $J_f$
pieces per signed grazing strip.
This reasoning uses actual first-hit inequalities for all earlier
disks, not an independently reflected comparison path.

Compactness of the radius interval and the outgoing collision
cylinder supplies a common neighborhood diameter $\eta$ from this
finite physical local argument.  It is compactness at one flight,
not compactness of pressure sums, invariant spaces or mixing times.
On the possible $H_0$ piece the adapted contraction is at most
$\Lambda^{-1}$; on $H_k^\pm$ it is at most
$(2a\tau_- k^2)^{-1}$, by the exact differential.
Summing all pieces, both signs, gives \eqref{eq:wp-cell-step}.
A vertical side maps to a positive cone direction
$\varkappa'+c'/\tau$; a horizontal side maps to
$\varkappa'+c'\varkappa/(\tau\varkappa+c)$.
The same complete-barrier argument therefore covers the indicated
rectangle restrictions and their descendants.
\end{proof}

\subsection{Matching the actual whole-cell comparison}

\begin{lemma}[Uniform comparison on the actual backward partition]
\label{lem:wp-distortion}
There is $C_{\rm cell}<\infty$, independent of $R,n$, such that
for two homogeneous short cone-stable curves $U,V$ in the same
$\mathcal M_{-n}^{0,H}$ cell, and regular points $u\in U,v'\in V$,
\[
 C_{\rm cell}^{-1}\le
       \frac{J_UT_R^{-n}(u)}{J_VT_R^{-n}(v')}
                       \le C_{\rm cell}.
\]
Curves may be different and may be restricted to arbitrary smaller
open pieces.  True singular points are excluded.
\end{lemma}

\begin{proof}[Primary-text matching and constant calculation to be completed]
This is the geometric content of Baladi--Demers, Lemma 2.2, applied
to its literal partition (2.4), not to an arbitrary physical cylinder.
Its fixed-table statement alone does not assert uniformity in $R$.
The proof on pages 7--8 uses the following three bounds:
\[
 d(T^{-i}u,T^{-i}v')
   \le C_b(\Lambda^{-i}+\Lambda^{-n+i}),\quad
 |\log(c_i/\widetilde c_i)|
       \le C_h d_i^{1/3},\quad
 |\tau_i-\widetilde\tau_i|\le C_\tau d_i^{1/2}.
\]
The first is its two-sided-cell geometric bound (2.5); its constant
must be explicitly matched to the common cone and chart geometry.
The second follows from the actual $q=2$ strip widths.
For the third use the physical entry formula
\[
 q'=q+(t_z-\sqrt{R^2-b_z^2})v
\]
on each of the finitely many possible next-center branches and
$|\sqrt A-\sqrt B|\le\sqrt{|A-B|}$.
The positive lower radius, bounded candidate distances and common
flight bounds give one $C_\tau$.

The inverse slope graph transform has contraction at most $47/48$,
as in module 313, and the finite part of its coefficients has the
same uniform square-root modulus.
Consequently the sum over $i$ of the slope differences is bounded
by a fixed multiple of
\[
 \frac{2K+\sum_i d_i^{1/2}}{1-47/48}.
\]
In the exact logarithmic Jacobian formula the two Euclidean
$\sqrt{1+V^2}$ factors telescope, and every remaining denominator
is bounded below by $\tau_-a$.
The logarithmic incidence, flight and slope sums are therefore
bounded uniformly, because
$\sum_i(\Lambda^{-i}+\Lambda^{-n+i})^\beta$ is bounded
for $\beta=1/3,1/2$.
Exponentiation gives $C_{\rm cell}$.
No additional potential, spectral theorem, invariant measure
identification or uniform mixing rate is used.
The matching of the geometric bound (2.5), with the actual connected
cells and their boundary convention, is the external applicability
step to be independently reconstructed before this lemma is used.
\end{proof}

\subsection{Long transverse cells, including every grazing strip}

Fix $\delta>0$ small enough that
$(2+v^{-1})\delta\le\eta/4$, and then $n_0\ge1$ so that
$(1+v)C_s\Lambda^{-n_0}\le\eta/4$.
Call $B\in\mathcal B_n$, $n\ge n_0$, long if $|I_B|\ge\delta$ and
short otherwise.
No physical event is changed by this division of a positive upper sum.

\begin{lemma}[Uniform weighted testing of all long cells]
\label{lem:wp-tests}
For every $n\ge n_0$,
\[
       \sum_{\substack{B\in\mathcal B_n\\ |I_B|\ge\delta}}A_n(B)
                          \le L_*<\infty
\]
with one constant independent of $R,n$.
Every homogeneity strip is included.
\end{lemma}

\begin{proof}
In each of the at most $N_f$ physical charts, choose a grid of
$h$-levels with spacing $\delta/3$.
The full range of $h$ has length at most $\pi(1+vR_+)$, so there
are at most
\[
 N_h=3+\left\lceil3\pi(1+vR_+)/\delta\right\rceil
\]
levels.  Intersect each level with every $H_k^\pm$ and with $H_0$,
and subdivide into overlapping short admissible stable curves.
The number per level and per strip is bounded by one common
integer $N_s$; for example use the chart-line length
$\pi R_+\sqrt{1+v^2}$ and the fixed admissible length.
The curves have zero curvature and slope $-v$.
At a fixed radius a generic translation of the grid makes $M_R'$
full arclength on every one of these countably many tests, by
Fubini applied to the affine stable foliation.
The translation preserves the mesh bound and is chosen once
before $n$.

A long connected cell has an open interval of $h$-values of length
at least $\delta$, so it meets one test in an open piece.
Assign it to its first such test, then to one homogeneous inverse
component $V\in\mathcal G_n(W)$.
Different backward cells assigned to the same test give different
inverse components: the complete inverse partition cuts every
backward true singularity and homogeneous boundary that separates
them.  Length subdivisions only refine this assignment.

At $y\in B\cap M_R'$ use a short affine curve with the true stable
slope at $y$, contained in $B$.  The preceding whole-cell comparison,
inverted, gives
\[
 j_n(T^{-n}y)\le
       C_{\rm cell}\sup_V J_VT^n.
\]
The exact identity \eqref{eq:wp-area} then shows
\[
 A_n(B)\le e_0^{-1}(\sup_{H_k}c)\,
                 C_{\rm cell}\sup_VJ_VT^n.
\]
Apply the existing stable inverse version of module 283 to each
test; its growth constant is $C_g$ for all short curves, including
arbitrarily short strip intersections.
Finally
\[
 C_H:=\sum_{H\in\mathcal H}\sup_H c
 \le1+2\sum_{k\ge k_0}k^{-2}<\infty.
\]
It follows that
\[
 L_*=e_0^{-1}C_{\rm cell}N_fN_hN_s C_g C_H
\]
is valid.  Dropping the terminal incidence before summing the strip
tests would invalidate this argument; it is retained throughout.
\end{proof}

\subsection{The complete short-cell renewal}

\begin{lemma}[A finite initial budget and summable short descendants]
\label{lem:wp-renewal}
There are $M_0<\infty$ and $K_0<\infty$, depending only on the
one-flight geometry and $\eta$, such that
\[
 S_l\le K_0^l C_H\quad(l\ge0)
\]
and, if $B\in\mathcal B_l$, $l\ge n_0$, then the descendants at time
$l+j$ having no long ancestor after $B$ have total relative inverse
unstable weight at most
\[
       \Gamma M_0\theta^j .
\]
If $B$ itself is short, $M_0$ is unnecessary.
\end{lemma}

\begin{proof}[Candidate proof; refined cell assignment to be checked]
Use a fixed cover by coordinate rectangles of diameter at most
$\eta$, with a finite multiplicity and at most $M_0$ rectangles
meeting any physical collision chart.
Intersect $B$ with these rectangles.
By the complete-barrier description, these are connected band
restrictions; their images have the barrier form in
Lemma~\ref{lem:wp-small}.
For each source restriction that is small, use the complete
one-step adapted inequality \eqref{eq:wp-cell-step}.
The adapted norm is between the Euclidean norm and $\Gamma$ times
that norm.  Converting only at the endpoints of a $j$-step chain
gives $\Gamma\theta^j$, never $\Gamma^j\theta^j$.

If each later original cell is short and its past length is at
least $n_0$, \eqref{eq:wp-diameter} makes it small.
Every refined descendant is contained in such an original cell.
Induction of the positive one-step sum thus proves the chain
bound; the initial rectangle restriction costs at most $M_0$.
If $B$ is already short, start the induction on $B$ itself.

For the finite initial budget no decay is needed.
The same finite rectangle subdivision at every step gives a
finite global one-step bound, for example
$K_0=\max\{1,\Gamma M_0\theta\}$.
The exact cocycle identity for $r_n$ and positivity give
$S_{l+1}\le K_0S_l$.
A coarse descendant may meet more than one rectangle image, but
its supremum is at most the sum over those refined intersections;
all intersections are counted on the upper side.
Since $S_0=C_H$, induction proves the displayed budget.
\end{proof}

\begin{theorem}[Conditional conclusion of the candidate geometry]
\label{thm:wp-conclusion}
If the independently reconstructed lemmas above hold with their
stated actual-cell scopes, then the first pressure assertion of
Input~\ref{input:v112-geometric-pressure} is proved:
\[
                   \sup_{R,m}P_m(R)\le Q_*<\infty.
\]
The theorem is unconditional only after the flagged geometric
proofs, rather than their conclusions, have been checked.
\end{theorem}

\begin{proof}
At time $n\ge n_0$ separate the long and short cells.
A short cell either has no long ancestor since time $n_0$, or a
unique most recent long ancestor at time $n-j$, $j\ge1$.
Its actual homogeneity and physical past determine every ancestor.
All refinements in Lemma~\ref{lem:wp-renewal} are positive upper
refinements and do not create a new orbit.

Using the exact cocycle identity, the last-long decomposition gives
\[
 S_n\le L_*+\Gamma\theta^{n-n_0}S_{n_0}
       +\Gamma M_0\sum_{j=1}^{n-n_0}
             \theta^j
             \sum_{\substack{B\in\mathcal B_{n-j}\\ B\ {\rm long}}}
                    A_{n-j}(B).
\]
Thus
\[
 S_n\le L_*+\Gamma K_0^{n_0}C_H
                 +\Gamma M_0 L_*\frac{\theta}{1-\theta}
                 =:S_*.
\]
The finitely many earlier $n$ satisfy the same type of bound by the
finite initial budget.  No positive fraction of current long cells,
mixing time or pressure limit is used.

For each physical word of length $n+1$, choose a point approaching
its stable-Jacobian supremum.
At its time-$n$ endpoint $y$ the past homogeneous cell determines
the first $n$ actual physical symbols; the last symbol has at most
$D$ choices.  Thus at most $D$ selected words are assigned to any
one backward cell.
The two optical/stable inequalities already proved in module 313
and \eqref{eq:wp-area} give
\[
 j_{n+1}(x)=j_n(x)j(T^nx)
       \le6400\,c(y)j_n(x)
       \le(6400/e_0)\,r_n(y).
\]
After the finite physical-word sum and then the supremum approximation,
\[
 P_{n+1}(R)\le(6400D/e_0)S_n(R).
\]
Combine this with $P_0=1$ and the finite initial bounds to obtain
$Q_*$.  The exact-label and conditioning tests have never been
changed: this is an unconditioned positive upper pressure estimate
used only where the original proof calls for it.
It does not itself prove the collision-normalized shortage limit,
the complete twisted raw local limit or the four-coordinate inverse
theorem.
\end{proof}

